\documentclass[10pt,a4paper]{amsart}
\usepackage[margin=19mm]{geometry}
\usepackage{amsmath,amssymb,mathtools}
\usepackage[scr=rsfs,cal=euler]{mathalfa}
\usepackage{xfrac}
\usepackage[unicode,hidelinks,bookmarks=false]{hyperref}
\newcommand{\ie}{i.e.\ }
\newcommand{\cf}{cf.\ }
\newcommand{\resp}{respectively\ }
\newcommand{\Subsection}{\subsection}
\providecommand{\nobreakdash}{\nobreak\hbox{-}}
\setlength{\emergencystretch}{2em}
\multlinegap0pt
\usepackage{algorithm}\usepackage{algpseudocode}
\usepackage{amssymb}
\usepackage{mathtools}
\usepackage{xcolor}
\usepackage{bm}
\newtheorem{proposition}{Proposition}
\newcommand{\z}{\mathbb{Z}}
\newcommand{\zz}{\mathbb{ZZ}}
\title{Quasi Zigzag Persistence: A Topological Framework for Analyzing Time-Varying Data}
\author{Tamal K. Dey, Shreyas N. Samaga}
\date{Source: arXiv:2502.16049v3 (CC BY 4.0)}
\begin{document}
\maketitle

\noindent\emph{Source and licence.} Quasi Zigzag Persistence: A Topological Framework for Analyzing Time-Varying Data. Version arXiv:2502.16049v3 (CC BY 4.0), distributed by arXiv under the Creative Commons Attribution 4.0 International licence, http://creativecommons.org/licenses/by/4.0/, as recorded in the arXiv licence metadata for this exact version and shown on the abstract page https://arxiv.org/abs/2502.16049v3. Full licence notice: \texttt{licenses/arxiv-2502.16049-CC-BY-4.0.txt}.\par\smallskip

\noindent\textit{Editorial scope.} Counted unit: the article's proposition aprop:init\_functor (source0.tex lines 683--686). The article's complete original proof of it is delivered with it (source0.tex lines 687--707, one \texttt{proof} environment of 21 lines). No mathematics is written, repaired or re-derived: the statement and the proof are the article's own lines, in the article's own notation, and no retained line is substituted or re-flowed.  No further source-local context is needed: every cross-reference of every retained line resolves inside the two delivered spans.  Nothing was omitted from any retained span, so the package carries no omission record.  No retained line contains a citation, so the wrapper carries no bibliography and no bibliography entry.  No retained line contains a figure environment, an image-inclusion command or drawing code, so no image is delivered and the package declares no asset dependency.  Editorial formatting only, in addition: an AMS \texttt{amsart} preamble that restates the article's own declarations and macros used by the retained lines, namely the environments \texttt{proposition} on separate counters, together with the standard packages those lines need.  The retained environments print in this document as Proposition 1 in order of appearance, whereas the article prints the same items with the numbering of its own full document.  The complete native source of the article as distributed in the arXiv e-print for this exact version is bundled unchanged as \texttt{sources/173699/source0.tex}.
\begin{proposition}
\label{aprop:init_functor}
    $F_L: \partial_L I\rightarrow I$  is an initial functor and $F_U: I\rightarrow \partial_U I$ is a terminal functor.
\end{proposition}

\begin{proof}
    It is sufficient to prove that $\partial_L I$ is initial for $F_L$ and $\partial_U I$ is terminal for $F_U$. We only show that $\partial_L I$ is initial and the proof
    for $\partial_U I$ being terminal is similar.

    Let $p\in I\setminus \partial_L I$ be any point. Let $p^-$ and $p^+$ be two points (if exist) where $p^-\prec p$ and $p^+ \succ p$ in $\z^2$ with $p^-_x < p_x < p^+_x$. By definition of the quasi zigzag poset, either $p^- > p$ and $p^+ > p$, or $p^- < p$ and
    $p^+ < p$ in the poset $\zz\times \z$.
    
    
    In the first
    case, the downset $\downarrow p$ consists of all points $p'\neq p$ that are vertically below $p$, that is, 
    $p'_x=p_x$ and $p'_y< p_y$. This means $\downarrow p$ intersects $\partial_L I$
    in a connected poset $p'_0 <\cdots < p'_k$ where each $p'_i$ has the same $x$-coordinate as $p$.

    In the second case, the downset $\downarrow p$ consists of three sets of points $Y^{p^-}$, $Y^p$, and $Y^{p^+}$  
    that are vertically below $p^-$, $p$, and $p^+$ respectively. Each of $Y^{p^-}$,
    $Y^p$ and $Y^{p^+}$ intersects $\partial_L I$ in a connected poset. They can be
    disconnected as a whole only if there is a point $q\in \partial_L I$ 
    so that
    $p^-_x < q_x < p_x$ or $p_x < q_x < p^+_x$. But, neither is possible because
    $p^-$ and $p^+$ are points in $I$ covered by $p$.    
\end{proof}

\end{document}
